Fix an even horizon \(T=2n\) and use the construction in the preceding proposition. To compare it with Saco's binary notation, define \(A_t^{\mathrm S}:=\mathbf1\{A_t=1\}\), write the executed propensity as \(\pi_t^{\mathrm S}:=p_t\), and relabel the fixed potential outcomes as \(Y_t^{\mathrm S}(1):=Y_t(1)=1\) and \(Y_t^{\mathrm S}(0):=Y_t(2)=1\). We verify the source hypotheses one by one.

First, \((X_t,Y_t^{\mathrm S}(0),Y_t^{\mathrm S}(1))\) is the same deterministic vector at every occasion. It is therefore independent and identically distributed, and the observed response satisfies \(Y_t=Y_t^{\mathrm S}(A_t^{\mathrm S})=1\); hence the source's superpopulation-arrival, consistency, and no-interference conditions hold. Its average treatment effect is
\[
 \theta_0=E[Y_t^{\mathrm S}(1)-Y_t^{\mathrm S}(0)]=1-1=0=\theta^*.
\tag{32}
\]

Second, \(p_t\) is \(\mathcal F_{t-1}\)-measurable and is the probability actually passed to the randomization device. If \(\mathcal G_t^{\mathrm S}:=\sigma(\mathcal F_{t-1},X_t,p_t)\) is Saco's pre-treatment information, then
\[
 P(A_t^{\mathrm S}=1\mid\mathcal G_t^{\mathrm S})=p_t=\pi_t^{\mathrm S}.
\tag{33}
\]
Because the potential outcomes are deterministic, \(A_t^{\mathrm S}\) is conditionally independent of them given \(\mathcal G_t^{\mathrm S}\). Thus logging integrity and sequential randomization hold. The policy uses only \(p_t\in\{1/2,1/4\}\), so
\[
 \frac14\le p_t\le\frac12,\qquad
 \frac12\le1-p_t\le\frac34.
\tag{34}
\]
This verifies executed overlap with \(\varepsilon=1/4\).

Third, choose the scored set \(\mathcal T_T:=\{1,\ldots,T\}\). It is deterministic, \(n_{\mathrm{eff}}=T\to\infty\), and no optional or data-dependent scoring is present. Choose
\[
 \widehat m_{t-1,0}(x)=\widehat m_{t-1,1}(x)=0
\tag{35}
\]
for every \(t,x\). These fits are deterministic and hence \(\mathcal F_{t-1}\)-measurable. Their fourth moments are zero, so both predictable nuisance construction and nuisance stability hold. Since \(|Y_t^{\mathrm S}(a)|^4=1\) for \(a\in\{0,1\}\), the source's potential-outcome fourth-moment condition also holds.

Fourth, substituting (32)--(35) into Saco's AIPW pseudo-outcome gives
\[
 \begin{aligned}
 \widehat\phi_t
 &=\frac{A_t^{\mathrm S}}{p_t}\{Y_t-0\}
   -\frac{1-A_t^{\mathrm S}}{1-p_t}\{Y_t-0\}+0-0\\
 &=\frac{\mathbf1\{A_t=1\}}{p_t}
   -\frac{\mathbf1\{A_t=2\}}{1-p_t}
 =\xi_t.
 \end{aligned}
\tag{36}
\]
Because \(\theta_0=0\), Saco's centered AIPW increment \(\widehat\phi_t-\theta_0\) is exactly the existing \(\xi_t\). In particular \(|\xi_t|\le4\), so \(\sup_tE|\xi_t|^4\le256\); this also verifies directly the bounded fourth-moment consequence used in the source proof.

By (36) and the definitions of the two quadratic variations,
\[
 S_{\mathcal T_T}^{\mathrm S}=\sum_{t=1}^T\xi_t=S_T,\qquad
 Q_{\mathcal T_T}^{\mathrm S}=\sum_{t=1}^T\xi_t^2=[S]_T,\qquad
 (V_{\mathcal T_T}^{\mathrm S})^2=\sum_{t=1}^TE[\xi_t^2\mid\mathcal F_{t-1}]=\langle S\rangle_T.
\tag{37}
\]
The preceding proposition proves pathwise that \(\langle S\rangle_T\ge4T\). Thus Saco's variance-growth assumption holds with \(v_-=4\) and failure probability zero. It also proves \(T^{-1}([S]_T-\langle S\rangle_T)\to_P0\). Since \(\langle S\rangle_T\ge4T\),
\[
 \left|\frac{[S]_T}{\langle S\rangle_T}-1\right|
 \le\frac{|[S]_T-\langle S\rangle_T|}{4T}\to_P0.
\tag{38}
\]
For completeness, with the Lyapunov exponent \(\delta=2\), the remaining ratio displayed in Saco's Appendix A is bounded pathwise by
\[
 \frac{\sum_{t=1}^TE|\xi_t|^4}{\langle S\rangle_T^2}
 \le\frac{256T}{16T^2}=\frac{16}{T}\longrightarrow0.
\tag{39}
\]
Thus the failure below is not caused by lack of overlap, moments, predictable scoring, quadratic-variation growth, realized-to-predictable quadratic-variation consistency, or the displayed Lyapunov bound.

The preceding proposition nevertheless proves, along \(T=2n\),
\[
 \frac{S_{\mathcal T_T}^{\mathrm S}}{\sqrt{Q_{\mathcal T_T}^{\mathrm S}}}\Longrightarrow L,\qquad
 E[L]=\frac2{\sqrt{2\pi}}\left(\sqrt{\frac3{28}}-\frac1{\sqrt8}\right)\ne0.
\tag{40}
\]
The limit \(L\) therefore does not have the \(N(0,1)\) law. Since convergence of the full sequence to \(N(0,1)\) would force the same limit along every subsequence, (40) contradicts the first conclusion asserted in Saco's Theorem 5.14.

The source's practical Studentized statistic fails as well. In its notation \(\widehat\theta^{\mathrm S}=T^{-1}S_T\), and its sample variance \(\widehat V^{\mathrm S}\) obeys the exact identity
\[
 (T-1)\widehat V^{\mathrm S}=Q_T-\frac{S_T^2}{T}.
\tag{41}
\]
Equation (40) implies \(S_T/\sqrt{Q_T}=O_P(1)\); consequently
\[
 \frac{T\widehat V^{\mathrm S}}{Q_T}
 =\frac{T}{T-1}\left\{1-\frac1T\left(\frac{S_T}{\sqrt{Q_T}}\right)^2\right\}\to_P1.
\tag{42}
\]
Slutsky's lemma, (40), and (42) give
\[
 \frac{\sqrt T(\widehat\theta^{\mathrm S}-\theta_0)}{\sqrt{\widehat V^{\mathrm S}}}
 =\frac{S_T/\sqrt{Q_T}}{\sqrt{T\widehat V^{\mathrm S}/Q_T}}\Longrightarrow L,
\tag{43}
\]
not \(N(0,1)\). This establishes the asserted exact in-class counterexample to the source theorem.